% T26 — «Морская тематика»: волны в шапке, нумерация в «спасательном круге».
% Для начальной школы. Math, класс 5, 6 простых задач.
\documentclass[a4paper,12pt]{article}

\usepackage{../../../../Lessons/_templates/preamble_worksheet}
\usepackage{../_shared/worksheet_check}

\usepackage{eso-pic}

% --- Палитра: морской синий ---
\definecolor{wcprimary}{HTML}{0B4F8A}
\definecolor{wcprimaryLight}{HTML}{8FCDEA}
\definecolor{wcprimaryBG}{HTML}{E8F4FB}
\definecolor{wcsand}{HTML}{F2E2B6}

\renewcommand{\familydefault}{\sfdefault}

% --- Светло-голубой фон страницы ---
\AddToShipoutPictureBG{%
  \AtPageLowerLeft{%
    \color{wcprimaryBG}\rule{\paperwidth}{\paperheight}}}

% --- Заголовок с волнами ---
\renewcommand{\WorksheetTitle}[1]{%
  \par\noindent\begin{tikzpicture}
    % Тёмная полоса-небо
    \fill[wcprimary] (0,0) rectangle (\linewidth, -18mm);
    % Волны (синусоида)
    \draw[wcprimaryLight, line width=1.4pt, domain=0:\linewidth, samples=120, smooth]
      plot (\x, {-15mm + 1.5mm*sin(\x*30)});
    \draw[white, line width=0.8pt, domain=0:\linewidth, samples=120, smooth]
      plot (\x, {-17mm + 1.2mm*sin(\x*30 + 60)});
    % Солнце-кружок справа
    \fill[wcsand] ([xshift=-12mm,yshift=-7mm]\linewidth,0) circle (4mm);
    \node[anchor=west, text=white, font=\bfseries\Large\sffamily]
      at (4mm, -6mm) {МОРСКАЯ МАТЕМАТИКА};
    \node[anchor=west, text=wcprimaryLight, font=\bfseries\large\sffamily]
      at (4mm, -12mm) {#1};
  \end{tikzpicture}\par\vspace{4mm}}

% --- TaskBox: номер в «спасательном круге» (concentric circles) ---
\renewcommand{\TaskBox}[2]{%
  \par\noindent\begin{tikzpicture}
    \node[draw=wcprimary, line width=1pt, rounded corners=6pt,
          fill=white, inner sep=4mm,
          text width=\dimexpr\linewidth-22mm\relax,
          anchor=north west] (B) at (12mm,0) {#2};
    % Спасательный круг
    \draw[wcprimary, line width=1.8pt, fill=white] (5mm, -5mm) circle (5mm);
    \draw[wcsand, line width=1.4pt] (5mm, -5mm) circle (3.5mm);
    \node[font=\bfseries\large\color{wcprimary}] at (5mm, -5mm) {#1};
  \end{tikzpicture}\par\vspace{3mm}}

\SetAnswerFieldSize{32mm}{9mm}

\begin{document}

\WorksheetTitle{Путешествие капитана Игрека}
{\small\itshape\color{wcprimary!85}Класс 5 \quad·\quad T26 \quad·\quad ID: T26-sea-2026-05-27}\par\vspace{2mm}
\FirstPageHeader

\TaskBox{1}{%
  Корабль вышел из порта в $7$ часов утра и прибыл к острову в $11$ часов того же дня.
  Сколько часов длилось плавание?\par
  \vspace{1.5mm}\hfill\textbf{\color{wcprimary}Ответ:}~\AnswerField{1}%
}

\TaskBox{2}{%
  В трюме корабля $48$ ящиков с апельсинами. Их разложили поровну на $6$ полок.
  Сколько ящиков на каждой полке?\par
  \vspace{1.5mm}\hfill\textbf{\color{wcprimary}Ответ:}~\AnswerField{2}%
}

\TaskBox{3}{%
  Рыбак поймал $23$ рыбы, а его сын~--- на $8$ рыб меньше.
  Сколько всего рыб у них в лодке?\par
  \vspace{1.5mm}\hfill\textbf{\color{wcprimary}Ответ:}~\AnswerField{3}%
}

\TaskBox{4}{%
  Маяк светит каждые $15$ секунд. Сколько раз он мигнёт за $3$ минуты?\par
  \vspace{1.5mm}\hfill\textbf{\color{wcprimary}Ответ:}~\AnswerField{4}%
}

\TaskBox{5}{%
  В порту стоят $7$ парусников и $12$ катеров.
  На сколько катеров больше, чем парусников?\par
  \vspace{1.5mm}\hfill\textbf{\color{wcprimary}Ответ:}~\AnswerField{5}%
}

\TaskBox{6}{%
  Капитан проплыл $128$ морских миль за $4$ дня поровну.
  Сколько миль он проходил за день?\par
  \vspace{1.5mm}\hfill\textbf{\color{wcprimary}Ответ:}~\AnswerField{6}%
}

\WorksheetQR{T26-sea/L1/v1}

\end{document}
